\documentclass[11pt]{article}
\usepackage[margin=0.75in]{geometry}
\usepackage{enumitem}
\usepackage{hyperref}
\usepackage{parskip}
\setlength{\parskip}{4pt}

\title{\vspace{-2em}Counterfactual Explainability: Deep RL vs.\ Online Planning\vspace{-1em}}
\author{Sean Florez (\texttt{sean.florez@colorado.edu})}
\date{}

\begin{document}
\maketitle
\thispagestyle{empty}

\section*{Problem Statement}

When a DQN agent picks an action, we know what it chose but not why.
The Q-values live inside neural network weights, so the reasoning is opaque.
Online planners like MCTS are different: the search tree stores simulated outcomes for every action considered, so the reasoning process is recorded and inspectable.
This project investigates a natural question: is suboptimal online planning more explainable than suboptimal deep RL, specifically because the ``training data'' is in the tree?

We will work in \textbf{MiniGrid-Dynamic-Obstacles-8x8}, a gridworld where an agent navigates around randomly moving obstacles to reach a goal.
This is formulated as an \textbf{MDP}: state includes agent position/orientation, obstacle positions, and goal location; actions are \{turn left, turn right, move forward\}; reaching the goal gives reward scaled by speed, while hitting an obstacle gives $-1$ and ends the episode.
The dynamic obstacles and partial observability ($7 \times 7$ egocentric view) create near-collision states where decisions genuinely need explaining.

Our approach: train both a DQN agent and an MCTS planner on this environment, then build a system that generates counterfactual explanations for each agent's decisions using Monte Carlo rollouts and translates those into natural language with an LLM.

\section*{Levels of Success}

\begin{enumerate}[leftmargin=*, label=\textbf{\arabic*.}, itemsep=2pt]

\item \textbf{Minimum Working Example.}
Train a DQN on MiniGrid-Dynamic-Obstacles-8x8 (Stable-Baselines3).
Build a counterfactual rollout framework: at each step of the trained policy, run 100+ MC rollouts per available action and log mean return, collision rate, success rate, and steps to termination.
Produce statistical summaries showing why the chosen action was preferred over alternatives.

\item \textbf{Main Approach.}
Implement an MCTS agent for the same environment and extract equivalent counterfactual statistics from the search tree (visit counts, child values, simulated collision/success rates).
Build an LLM pipeline that converts rollout logs (DQN) and tree logs (MCTS) into natural language explanations.
Compare the two on explanation quality using fidelity and soundness metrics, following recent XRL evaluation work (Belouadah et al., 2025).
This directly tests whether tree-based planning produces better explanations than post-hoc analysis of a learned policy.

\item \textbf{Stretch Goal.}
Interactive chatbot where a user steps through a trajectory and asks questions like ``why did you turn left?'' or ``what if you went forward?''
Answers are grounded in rollout data and support follow-up questions for conversational exploration of agent behavior.

\end{enumerate}

\section*{Implementation Plan}

\textbf{From scratch:} counterfactual MC rollout framework and logging, MCTS agent for MiniGrid (building on HW3 MCTS work), LLM explanation pipeline and prompt design, explanation evaluation metrics, chatbot interface (stretch).

\textbf{Off-the-shelf:} MiniGrid environment (Farama Foundation), DQN training (Stable-Baselines3 or CleanRL), LLM API (Claude), standard Python stack (numpy, matplotlib, pandas).

\end{document}
